La glissade doit absolument passer par 4 points fixes ($A$,$B$,$C$,$D$) et un
point dont la hauteur est variable entre 10 et 15 mètres ($E$). Si les
coordonées exactes du point $E$ étaient connues à l'avance, créer un fonction
polynomiale passant par les 5 points pourrait être résolue en utilisant les
coordonées pour développer un système d'équations avec 5 équations et 5 inconnus
(polynôme d'ordre 4). Plusieurs suppositions sont émises dans MATLAB, chacunes
avec une valeur de $E$ différente pour qu'elles soient uniformément réparties
dans la plage permise. Pour chaque système d'équation, les coefficients du
polynôme sont résolus. Afin d'être bref seul la résolution pour le $E$ choisi
est montrée puisque la démarche est toujours la même (le choix sera justifié
plus tard):

\begin{align}
	F(x) &= a_0 + a_1 x + a_2 x^2 + a_3 x^3 + a_4 x^4\\
\begin{bmatrix} 30\\ 20\\ 20\\ 16\\ 12.5 \end{bmatrix}
&=
\begin{bmatrix}
1 & 0 & 0^2 & 0^3 & 0^4\\
1 & 7 & 7^2 & 7^3 & 7^4\\
1 & 15 & 15^2 & 15^3 & 15^4\\
1 & 20 & 20^2 & 20^3 & 20^4\\
1 & 25 & 25^2 & 25^3 & 25^4
\end{bmatrix}
\begin{bmatrix} a_0\\ a_1\\ a_2\\ a_3\\ a_4 \end{bmatrix} \Rightarrow
\begin{bmatrix} a_0\\ a_1\\ a_2\\ a_3\\ a_4 \end{bmatrix} =
\begin{bmatrix} 30.000000\\ -3.910623\\ 0.525220\\ -0.027685\\ 0.000473
\end{bmatrix} \\
	F(x) &= 30.000000 -3.910623 x + 0.525220 x^2 -0.027685 x^3 + 0.000473 x^4
\end{align}

Plusieurs équations polynomiales sont donc trouvées selon la hauteur du point
$E$ (Annexe \ref{fig:01-trajectories}). Toutefois, il est requis que la pente
finale soit horizontale ou presque. Cela correspond à la fonction polynomiale
pour laquelle la dérivée est la plus près de 0 à $x=25$. Ici, cela correspond à
$E_y=12.5$m (pente de $-0.0260$).

Puisqu'il est question ici d'une interpolation et non d'une régression, la
courbe passe par tout les points et il n'y a pas place pour des erreur. Autrement
dit, le RMS est de 0 par définition.\\

